\chapter*{Formelzeichen und Abk\"urzungen}

\section*{Formelzeichenverzeichnis}
\enlargethispage{-4cm}

\renewcommand{\arraystretch}{1.3}
\setlength{\tabcolsep}{12pt}
\begin{longtable}[l]{l l p{8.5cm}}
\bfseries \emph{Formelzeichen}& \bfseries \emph{phys.Einheit} &\bfseries \emph{Bedeutung}\\
\endfirsthead
\bfseries \emph{Formelzeichen}& \bfseries \emph{phys.Einheit} &\bfseries \emph{Bedeutung}\\
\endhead
$a$	& $m$	&Seitenlänge\\
$A$ & $m^2$ & Fläche\\
$b$ & $m$ & Seitenlänge\\
$B$ & $N \cdot mm^2$ & Biegesteifigkeit\\
$c$ & $\frac{\text{N} \cdot \text{s}}{\text{m}}$ & mech. Dämpfung\\
$C$ & $.$ & Hilfsgröße $C = \sqrt{B/m''}$\\
$d$ & $m$ & Gesamtstärke der Platte\\
$E$ & $Pa$ & Elastizitätsmodul\\
$f$ & $Hz$ & Frequenz\\
$f_0$ & $Hz$ & Grundfrequenz\\
$f_{m,n}$ & $Hz$ & Eigenfrequenz der $(m,n)$-ten Mode\\
$f_s$ & $Hz$ & Abtastfrequenz\\
$\Delta f$ & $Hz$ & Frequenzauflösung des Spektrums\\
$F$ & $N$&	Kraft\\
$H(f)$ & $-$ & Übertragungsfunktion\\

$m$	& $kg$&	Masse\\
$m''$ & $\nicefrac{kg}{m^2}$ & Flächenmasse\\
$m, n$ & $-$ & Modenindizes\\

$n(f)$ & $\nicefrac{1}{Hz}$ & Modale Dichte\\
$p$ & $Pa$ & Akustischer Schalldruck\\
$p_0$ & $Pa$ & Bezugsdruck ($p_0 = 20 \cdot 10^{-6}$ Pa)\\

$r$ & $-$ & Seitenverhältnis der Platte\\

$t$ & $s$ & Zeit\\
$T$ & $mm$ & Stärke der Deckschicht\\

$\eta$ & $-$ & Verlustfaktor (Materialdämpfung)\\
$\sigma$ & $-$ & Abstrahlgrad (Abstrahleffizienz)\\

$k$ & $-$ & Frequenzindex\\
$N$ & $-$ & Anzahl der Abtastwerte pro Datenblock\\
$M$ & $-$ & Anzahl der Mittelungsblöcke\\
$X(k)$ & $-$ & Komplexes Spektrum des Referenzsignals\\
$Y(k)$ & $-$ & Komplexes Spektrum des Antwortsignals\\
$X_i(k)$ & $-$ & Komplexes Spektrum des Referenzsignals im $i$-ten Block\\
$X_i^*(k)$ & $-$ & Konjugiert komplexes Spektrum des Referenzsignals im $i$-ten Block\\
$Y_i(k)$ & $-$ & Komplexes Spektrum des Antwortsignals im $i$-ten Block\\
$A_i(k)$ & $Pa$ & Lineares Amplitudenspektrum des $i$-ten Datenblocks\\
$\bar{A}(k)$ & $Pa$ & Gemitteltes lineares Amplitudenspektrum\\
$G_{xx}(k)$ & $-$ & Autospektrum des Referenzsignals\\
$G_{yy}(k)$ & $-$ & Autospektrum des Antwortsignals\\
$G_{xy}(k)$ & $-$ & Kreuzspektrum zwischen Referenz- und Antwortsignal\\
$H_1(k)$ & $-$ & $H_1$-Schätzer der Übertragungsfunktion\\
$H_{\mathrm{dB}}(k)$ & $dB$ & Übertragungsfunktion in Dezibel\\
$L_p(k)$ & $dB$ & Schalldruckpegel (SPL)\\
$\gamma^2(k)$ & $-$ & Kohärenzfunktion\\

\end{longtable}

\clearpage
\section*{Abk\"urzungsverzeichnis}


\begin{table}[!h]
\setlength{\tabcolsep}{50pt}
	\centering
		\begin{tabular}[t]{l l}
\bfseries \emph{Abk\"urzung}& \bfseries \emph{Bedeutung}\\
		CAD	& Computer Aided Design\\
		DAQ	& Datenerfassungssystem\\
		DML	& Distributed Mode Loudspeaker\\
		FEM	& Finite-Elemente-Methode\\
		FFT	& Fast Fourier Transformation\\
		HiFi	& High Fidelity\\
		PML	& Perfectly Matched Layer\\
		RMS	& Root Mean Square\\
		SPL	& Sound Pressure Level\\
		VDI	& Verein Deutscher Ingenieure\\
		VI	& Virtuelles Instrument\\

		\end{tabular}

\end{table}
\normalsize
\normalfont
